\documentclass[10pt]{article}

\usepackage[utf8]{inputenc}

\usepackage[T1]{fontenc}

\usepackage{amsmath}

\usepackage{amsfonts}

\usepackage{amssymb}

\usepackage[version=4]{mhchem}

\usepackage{stmaryrd}

\usepackage{bbold}



\begin{document}

$\delta_{h}(x)=1$ при $x=0$ и $\delta_{h}(x)=0$ при $x=n h, n \neq 0$; а также решение задачи об излучении движущегося источника (см. также Осцилляций шлейф, Черенкова эффект в разностных уравнениях). Изложенные методы переносятся на более общие дифференциально-разностные, разностные и псевдодифференциальные уравнения. Они позволяют, напр., изучать устойчивость разностных схем в $C$.



Лит.: [1] Фе до р ю к М. В., «Журнал вычислительной математики и математич. физики», 1967, т. 7, № 3, с. 510-40; [2] М асло в В. П., «Успехи математич. наук», 1968 , т. 23 , в. 4, с. 243 ; [3] е го же, Операторные методы, М., 1973; [4] Дан и лов В. Г., Операторные методы исследования разностных схем, М., 1979; [5] Маслов В. П., Данилов В. Г., “Труды Математич. ин-та AH СССР», 1984, т. 166, c. 130-60; 1985, т. 167, c. 96-107; [6] Данилов В. Г., Жевандров П.Н., «Известия АН СССР. Сер. математическая», 1989, т. 53, № 2, с. 411-24.



В. Г. Данилов, П. Н. Жевандров.



PAЗPEШЕНИЕ ОСОБенНОСтей р остка $f:\left(\mathbb{R}^{n}, 0\right) \rightarrow$ $\rightarrow(\mathbb{R}, 0)$ аналитической функции в окрестности критической точки $O-$ это $n$-мерное многообразие $Y$ и аналитическое отображение $\pi: Y \rightarrow \mathbb{R}^{n}$ со свойствами:



\begin{enumerate}

  \item в подходящих координатах в окрестности любой точки из $\pi^{-1}(0)$ функция $f \circ \pi$ и якобиан $\pi$ являются одночленами, умноженными на не обращающуюся в 0 функцию;



  \item в окрестности точки 0 найдется собственное алгебраич. подмножество, вне к-рого $\pi$ аналитически обратимо;



  \item прообраз любого компакта в окрестности $x$ компактен.



\end{enumerate}



РАЗРЕШИМАЯ АЛГЕБРА ЛИ - см. Ли алгебра.



PАЗРЕШИМАЯ ГРУПпА - гpуnпа, в которой существует цепочка подгрупп



\[

G=G_{0} \supset G_{1} \supset \ldots \supset G_{k}=\{e\}

\]



таких, что $G_{i+1}$ является нормальной подгруппой в $G_{i}$ и факторгруппы $G_{i} / G_{i+1}$ абелевы. Для связных групп Ли обычно используется следующее эквивалентное определение: $G$ р аз реш и ма, если в ней существует цепочка замкнутых нормальных подгрупп (*) таких, что факторгруппы $G_{i} / G_{i+1}$ абелевы (через $е$ обозначена единица группы $G$ ). Группа Ли разрешима тогда и только тогда, когда ее Ли алгебра разрешима. Любая подгруппа Р. г. Ли разрешима. Любое неприводимое конечномерное представление Р. г. Ли одномерно. Любое конечномерное представление Р. г. Ли в подходящем базисе задается верхнетреугольными матрицами.



PA3PЫBA ПОВЕРХНОСТИ - поверхности, на которыХ терпят разрыв параметры среды. Р. п. возникают, в частности, в нелинейной звуковой волне вследствие того, что точки, соответствующие большей плотности, движутся быстрее, чем точки, соответствующие меньшей плотности. Благодаря этому профиль волн меняет со временем свою форму. Точки сжатия выдвигаются вперед, а точки разрежения оказываются отставшими. В конце концов профиль волны изгибается настолько, что кривая $\rho(x)$ (в заданный момент времени) оказывается неоднозначной: нек-рым $x$ соответствуют три различных значения $\rho$. Такая неоднозначность невозможна: в этом случае функция $\rho(x)$ становится разрывной, в результате чего становится однозначной функцией. Р. п. возникают также вследствие разрыва в начальных условиях, напр. при столкновении двух газовых масс.



На Р.п. дифференциальные уравнения теряют смысл и решения по обе стороны Р. п. сшиваются с помощью граничных условий. В обычной гидродинамике существуют два типа Р. п.- тангенци альные, через поверхность к-рых нет потока вещества, и ударные волны, у к-рых частицы вещества пересекают Р. П. При формулировке граничных условий на ударных волнах исходят из закона сохранения энтропии, полагая газ идеальным. В действительности на поверхности ударной волны происходит диссипация превращение части кинетич. энергии газа в тепло. Поэтому на ударной волне сохраняется не энтропия, а энергия.



Количественно Р. П. описываются скачками параметров среды - разностыю их значений по обе стороны Р. п. Вели- чины этих скачков определяются законами сохранения массы, импульса и энергии на Р. п.



При учете диссипации (вязкости и теплопроводности) никаких разрывов в принципе не может быть и все течения непрерывны. Однако при малой диссипации происходит очень быстрый (в узком слое) непрерывный переход от одного состояния к другому, к-рый приближенно рассматривается как разрыв.



Коэффициенты диссипации входят в уравнения гидродинамики множителями при производных гидродинамич. величин. На Р. П. эти производные обращаются в бесконечность. Поэтому диссипативные члены не малы, а представляют собой неопределенность типа $0 \cdot \infty$, то есть ими нельзя пренебрегать.



Лит.: [1] Ландау Л. Д., Л ифшиц Е. М., Гидродинамика, 3 изд., М., 1986; [2] С едов Л. И., Механика сплошной среды, т. 1-2, 4 изд., М., 1983-84; [3] Курант Р., Фридрихс К., Сверхзвуковые течения и ударные волны, пер. с англ., М., 1950.



В. П. Демуцкий, Р. В. Половин.



PАЗРЫВНОЕ КОЛЕБАНИЕ - см. Релаксационное колебание.



PАЗРЫВНЫЙ ПРОЦЕСС - см. Саучайный процесс.



РАЙФЕРТИ ТЕОРЕМА - см. Внутренняя симметрия.



PAHг симметрического пространства - см. Симметрическое пространство.



PAнг те нзо ра - см. Тензорное поле.



PACKPЫTЫй 30нтИК - особая двумерная поверхность $\Lambda^{2}$ в $T^{*} \mathbb{R}^{2} \cong \mathbb{R}^{4}$, являющаяся пространством конормального расслоения полукубической параболы $\left\{x^{2}=y^{3}\right\} \subset \mathbb{R}^{2}$, то есть множество ковекторов, приложенных к точкам полукубической параболы и ортогональных ей. Образ Р. з. при проектировании $\mathbb{R}^{1} \rightarrow \mathbb{R}^{3},\left(x, y, p_{x}, p_{y}\right) \mapsto\left(y, p_{x}, p_{y}\right)$ является зонтиком Уитни $4 p_{x}^{2}=9 y p_{y}^{2}$. Обобщенный Р.з. определяется как пространство конормального расслоения (обобщенного) раскрытого ласточкина хвоста. Р. з. является важным примером особого лагранжева многообразия и появляется в качестве особенности систем лучей.



Пусть $H-$ неособая гиперповерхность в симплектич. многообразии $\left(M^{2 n}, \omega\right), \quad l \subset H-$ гладкое изотропное $(n-1)$-мерное подмногообразие общего положения. Тогда многообразие характеристик поверхности $H$, проходящих через $l$ - лагранжево многообразие, особенности к-рого симплектически изоморфны особенностям (обобщенных) $\begin{array}{ll}\text { Р. з. или их надстроек. } & \text { М. Э. Казарян. }\end{array}$



РАСКРЫТЫЙ ЛАСТОЧКИН ХВОСТ - особое мНогообразие $\Sigma^{k}$ в пространстве многочленов



\[

x^{2 k+1}+a_{2 k-1} x^{2 k-1}+\ldots+a_{1} x+a_{0}

\]



состоящее из многочленов, имеющих $(k+1)$-кратный корень. Р.л. Х. -лагранжево многообразие относительно нек-рой канонически определенной симплектической структуры в пространстве многочленов. Многообразие многочленов, имеющих $k$ пар кратных корней, также лагранжево (относительно нек-рой др. симплектич. структуры) и симплектически изоморфно Р. л. х. Многообразие $\Sigma^{1}-$ полукубич. парабола на плоскости. Р. л. х. является важным примером особого лагранжева многообразия и появляется в качестве особенности систем лучей в задаче об обходе препятствий.



Пусть $H$ - неособая гиперповерхность в симплектич. многообразии $\left(M^{2 n}, \omega\right), L \subset M$ - лагранжево многообразие, квадратично касающееся $H$ вдоль гладкого изотропного подмногообразия $l^{n-1}=L \cap H$, а в остальном общего положения. Тогда многообразие характеристик $H$, проходящих через $l$,лагранжево многообразие, особенности к-рого симплектически изоморфны особенностям Р. л. Х. или их надстроек.



Лum.: [1] Итоги науки и техники. Современные проблемы математики. Новейшие достижения, т. 33, М., $1988 . \quad$ М. Э. Казарян.





\end{document}